\documentclass[10pt]{IEEEtran}
\usepackage[utf8]{inputenc}
\usepackage{csquotes}
\usepackage[portuguese]{babel}
\usepackage{fontenc}

\usepackage[style=abnt-numeric, backend=biber, sorting=none]{biblatex}
\addbibresource{referencias.bib}
\DeclareFieldFormat{labelnumberwidth}{\mkbibbrackets{#1}}

% --- Classificação da Bibliografia ---

% 1. BASICA
\DeclareBibliographyCategory{basica}
\addtocategory{basica}{
    McKeown:OpenFlow,
    Kreutz:Comprehensive,
    Astuto:SDNSurvey,
    ONF_OpenFlow_1.5.1,
    Kurose:Redes,
    Tanenbaum:Redes
}

% 2. HARDWARE
\DeclareBibliographyCategory{hardware}
\addtocategory{hardware}{
    Forencich:Corundum,
    Zhong:Lightning,
    Saleh:AnalogINC,
    Mashreghi:PrismParser,
    10.1145/1477942.1477944
}

% 3. TENDENCIAS
\DeclareBibliographyCategory{tendencias}
\addtocategory{tendencias}{
    Bosshart:P4,
    Liatifis:P4Survey,
    Janabi:Security2024,
    Wang:TrustworthyEI,
    9972067
}

% 4. CONTEXTO
\DeclareBibliographyCategory{contexto}
\addtocategory{contexto}{
    Feamster:RoadToSDN,
    Casado:Ethane,
    Rothenberg:OpenFlow,
    Rothenberg:MininetWiFi,
    Guedes:SBRC2012,
    Das:10,
    6587999,
    Jammal:SDNChallenges,
    Kirkpatrick2013,
    IPJ_16_1
}

\title{Proposta de Referencial Teórico: SDN e Arquiteturas Pós-Digitais}
\author{
    \IEEEauthorblockN{Nico I. G. Ramos}
    \IEEEauthorblockA{
        Orientando - GRR20210574\\
    }
}
\date{\today}

\begin{document}
\maketitle

\section{Objetivo}
Este documento apresenta a seleção bibliográfica para a fundamentação teórica do Trabalho de Conclusão de Curso. A bibliografia está organizada em quatro categorias principais: (i) obras canônicas sobre SDN e OpenFlow; (ii) referências sobre aceleração de hardware e computação pós-digital; (iii) estudos sobre a evolução do paradigma, incluindo P4 e tendências emergentes; e (iv) contextualização histórica e regional do tema.

\section{Bibliografia Básica (Obras Canônicas)}
A seguir, detalha-se a justificativa acadêmica para a inclusão das referências estruturantes que definem o paradigma SDN e o protocolo OpenFlow.

\begin{itemize}
    \item \cite{McKeown:OpenFlow} \textbf{O Documento Fundador:} \\
    O trabalho seminal de McKeown et al. (2008) introduziu o OpenFlow como solução para a "ossificação" da Internet. Esta referência será a base para descrever a separação fundamental entre os planos de controle e dados.
    
    \item \cite{Kreutz:Comprehensive} \textbf{Survey de Referência:} \\
    A revisão extensiva de Kreutz et al. (2015) é essencial para discutir os requisitos de escalabilidade, confiabilidade e segurança na arquitetura SDN. A obra consolida o estado da arte da primeira década do paradigma.
    
    \item \cite{Kurose:Redes, Tanenbaum:Redes} \textbf{Livros Didáticos:} \\
    As obras de Kurose e Tanenbaum fornecem as bases teóricas de redes de computadores. Serão utilizadas pontualmente para garantir o rigor terminológico e a contextualização acadêmica dos conceitos de comutação e roteamento.
    
    \item \cite{Astuto:SDNSurvey, ONF_OpenFlow_1.5.1} \textbf{Definições e Especificações:} \\
    O \textit{survey} de Astuto et al. e a especificação oficial da ONF servirão como fontes primárias para a taxonomia padrão e para os detalhes técnicos das tabelas de fluxo e mensagens do protocolo.
\end{itemize}

\section{Aceleração de Hardware e Computação Pós-Digital}
Esta seção constitui o foco da atualização bibliográfica, selecionando obras fundamentais para o desenvolvimento de planos de dados de alto desempenho baseados em hardware reconfigurável e físico.

\begin{itemize}
    \item \cite{Forencich:Corundum} \textbf{FPGA e Open Hardware (Corundum):} \\
    O trabalho de Forencich et al. (2020) apresenta o \textit{Corundum}, uma NIC baseada em FPGA de código aberto operando a 100 Gbps. Esta obra será a referência central para a implementação de SDN em hardware moderno, demonstrando a superação das limitações das primeiras gerações (como o NetFPGA).
    
    \item \cite{Saleh:AnalogINC} \textbf{Computação Analógica na Rede:} \\
    Saleh et al. (2024) propõem o uso de \textbf{Memristores} e memórias associativas probabilísticas (pCAM) para processamento de pacotes no domínio analógico. Esta referência será utilizada para discutir a redução do consumo energético e as arquiteturas pós-digitais.
    
    \item \cite{Zhong:Lightning} \textbf{SmartNICs Fotônicas:} \\
    O projeto \textit{Lightning} (Zhong et al., 2023) demonstra uma SmartNIC híbrida que utiliza computação fotônica para inferência em tempo real. A obra exemplifica como eliminar gargalos eletrônicos no processamento de pacotes.

    \item \cite{Mashreghi:PrismParser} \textbf{P4 em FPGA:} \\
    Mashreghi-Moghadam et al. (2024) detalham arquiteturas eficientes para implementar \textit{parsers} P4 em FPGAs. Este trabalho é essencial para unir a flexibilidade da programação SDN com a velocidade de execução do hardware dedicado.
\end{itemize}

\section{Evolução do Paradigma (P4 e Tendências)}
Referências selecionadas para cobrir a programabilidade avançada do plano de dados e o estado da arte em segurança e gerenciamento.

\begin{itemize}
    \item \cite{Bosshart:P4} \textbf{Linguagem P4:} \\
    A proposta original de Bosshart et al. (2014) define os processadores de pacotes independentes de protocolo. Esta obra será utilizada para explicar a transição do modelo OpenFlow fixo para o modelo programável P4.
    
    \item \cite{Liatifis:P4Survey, Janabi:Security2024, Wang:TrustworthyEI} \textbf{Surveys Recentes e Segurança:} \\
    As revisões de Liatifis, Janabi e Wang (2023-2025) mapeiam a integração de SDN com IA e redes 6G. Serão utilizadas para demonstrar a atualidade da pesquisa e os desafios emergentes de segurança.

    \item \cite{9972067} \textbf{Gerenciamento de Filas em P4:} \\
    Com o avanço da programabilidade no plano de dados, novas aplicações tornaram-se viáveis diretamente nos switches. Kudure et al. (2022) exploram a implementação de algoritmos de Gerenciamento Ativo de Filas (AQM), como o CoDel, inteiramente em hardware P4, evidenciando a granularidade de controle atual.
\end{itemize}

\section{Contextualização Histórica e Regional}
Obras que situam o SDN na história intelectual das redes, destacam a produção científica brasileira e analisam os desafios de adoção.

\begin{itemize}
    \item \cite{Feamster:RoadToSDN} \textbf{História Intelectual:} \\
    Feamster et al. analisam a evolução das "Redes Ativas" até o SDN. Esta obra será utilizada na introdução para traçar a linha do tempo das tecnologias de gerenciamento de rede.
    
    \item \cite{Guedes:SBRC2012, Rothenberg:MininetWiFi} \textbf{Cenário Brasileiro:} \\
    Os trabalhos do SBRC e o desenvolvimento do Mininet-WiFi por Rothenberg et al. são fundamentais para contextualizar a contribuição nacional e descrever as ferramentas de emulação adotadas.
    
    \item \cite{Casado:Ethane, Rothenberg:OpenFlow} \textbf{Precursores e Motivação:} \\
    O projeto \textit{Ethane} (Casado et al., 2007) e a análise de Rothenberg (2013) justificam as motivações originais para a quebra de paradigma nas redes corporativas e acadêmicas.
    
    \item \cite{Jammal:SDNChallenges, Kirkpatrick2013, IPJ_16_1, 6587999} \textbf{Consolidação e Desafios:} \\
    Publicações seminais, como as apresentadas no \textit{Internet Protocol Journal} e por Kirkpatrick, destacam como o SDN alterou o paradigma de gerenciamento. A \textit{Open Networking Foundation} (ONF) reforçou essa visão estabelecendo padrões, enquanto Jammal et al. detalham os desafios de pesquisa que persistem na implementação dessas redes.

    \item \cite{Das:10} \textbf{Arquiteturas Unificadas:} \\
    Das et al. propuseram, ainda nos estágios iniciais, uma arquitetura de controle unificado capaz de gerenciar tanto redes de pacotes quanto de circuitos, demonstrando a flexibilidade inerente ao modelo SDN.
    
    \item \cite{10.1145/1477942.1477944} \textbf{Hardware Clássico:} \\
    O trabalho original sobre NetFPGA (Naous et al., 2008) é mantido aqui como referência histórica, servindo de contraponto para demonstrar a evolução do hardware discutido na Seção III.
\end{itemize}

\section{Listagem de Referências}
\nocite{*} 

% Imprime as bibliografias separadas por categoria para organização visual
\printbibliography[category=basica, title={Bibliografia Básica (Obras Canônicas)}]
\printbibliography[category=hardware, title={Hardware Avançado: FPGA, Analógico e Fotônica}]
\printbibliography[category=tendencias, title={Tendências: P4 e Evolução de Software}]
\printbibliography[category=contexto, title={Contexto: Histórico e Regional}]

\end{document}